\documentclass[UTF8,11pt,a4paper]{ctexart}
\usepackage{amsmath,amssymb,geometry,enumitem,booktabs,hyperref}
\geometry{margin=2.1cm}\setlist{itemsep=2pt,topsep=3pt}
\hypersetup{colorlinks=true,linkcolor=blue}
\newcommand{\EN}[1]{\par\noindent \textbf{English.} #1\par}
\newcommand{\READ}[1]{\par\noindent \textbf{公式英文读法 / How to read it.} #1\par}
\newcommand{\EX}[2]{\par\noindent \textbf{例子 / Example.} #1\par\EN{#2}}
\newcommand{\ASK}[4]{\par\noindent \textbf{可能问题 / Possible question.} #1\par\EN{#2}\noindent \textbf{回答要点 / Key answer.} #3\par\EN{#4}\medskip}
\title{拓扑学 APS 面试准备\\[4pt]\large 覆盖 SEU《Topology》全书范围}
\author{Topology: Bilingual Review}
\date{}
\begin{document}\maketitle

\section*{课程概述 | Course overview}
\noindent\textbf{中文口述稿。} 我在拓扑学中研究哪些性质在连续变形下保持不变。第一部分从度量空间推广到拓扑空间：开集定义了连续性，紧性可以保证极值存在，连通性描述空间是否能分成互不相交的两块。第二部分学习商拓扑和曲面，把多边形的边按规则粘合起来，例如把正方形相对边同向粘合得到环面，反向粘合可以得到 Klein 瓶或射影平面。最后用同伦、CW 复形和 Euler 特征区分曲面。Euler 特征是 $\chi=V-E+F$；球面为 $2$，环面为 $0$。这些例子让我理解拓扑学如何把“距离不重要、连接方式重要”的几何问题转成可计算的不变量。
\EN{In topology, I studied which properties remain unchanged under continuous deformation. The first part generalizes metric spaces to topological spaces: open sets define continuity, compactness guarantees extrema for continuous functions, and connectedness describes whether a space can be split into two separated pieces. The second part studies quotient spaces and surfaces by gluing edges of polygons. Identifying opposite sides of a square in the same direction gives a torus, while reversing an identification can give a Klein bottle or a projective plane. Finally, homotopy, CW complexes, and Euler characteristic distinguish surfaces. The Euler characteristic is $\chi=V-E+F$: it is two for a sphere and zero for a torus.}

\section{度量空间与拓扑空间 | Metric and topological spaces}
\noindent 度量空间 $(X,d)$ 中，$d$ 满足非负性、分离性、对称性和三角不等式。开球 $B(x,r)=\{y:d(x,y)<r\}$。拓扑空间 $(X,\mathcal T)$ 则直接指定哪些集合是开集：$\varnothing,X$ 开，任意并开，有限交开。度量会生成拓扑，但一个拓扑可能来自不同度量，也可能根本不能由度量产生。
\EN{A metric space has a distance satisfying positivity, separation, symmetry, and the triangle inequality. A topology directly specifies its open sets: the empty set and the whole space are open, arbitrary unions of open sets are open, and finite intersections are open. A metric induces a topology, but different metrics may induce the same topology and some topologies are not metrizable.}
\EX{实数的通常拓扑由开区间生成；离散拓扑中每个子集都开；平凡拓扑只有 $\varnothing$ 和 $X$。在离散拓扑中，序列收敛到 $x$ 当且仅当它最终恒等于 $x$。}{The usual topology on the real line is generated by open intervals. In the discrete topology every subset is open; in the trivial topology only the empty set and the whole space are open. In the discrete topology a sequence converges to $x$ exactly when it is eventually equal to $x$.}

\noindent\textbf{基。} 拓扑基 $\mathcal B$ 是一族开集，使每个点属于某个基元素，且两个基元素的交在每个交点附近还包含一个基元素。拓扑中的开集就是基元素的任意并。实数的开区间、度量空间的开球都是典型基。
\EN{A basis is a family of sets that covers the space and locally refines intersections. The open sets are arbitrary unions of basis elements. Open intervals in the real line and open balls in a metric space are standard bases.}

\section{闭集、连续性与子空间 | Closed sets and continuity}
\noindent 闭集是补集开集的集合；闭包 $\overline A$ 是包含 $A$ 的最小闭集，边界为 $\partial A=\overline A\setminus A^\circ$。映射 $f:X\to Y$ 连续，当且仅当任意开集 $V\subseteq Y$ 的原像 $f^{-1}(V)$ 在 $X$ 中开。等价地，任意闭集的原像闭。
\EN{A closed set has an open complement. The closure is the smallest closed set containing $A$, and the boundary is $\partial A=\overline A\setminus A^\circ$. A map is continuous exactly when inverse images of open sets are open, equivalently when inverse images of closed sets are closed.}
\EX{函数 $f(x)=x^2$ 从 $\mathbb R$ 到 $\mathbb R$ 连续，因为开集的原像开；但它不是一一的，$f(1)=f(-1)$。连续性不等于同胚，双射和逆映射连续还要另外检查。}{The map $x^2:\mathbb R\to\mathbb R$ is continuous but not injective, since $f(1)=f(-1)$. Continuity alone is not a homeomorphism; bijectivity and continuity of the inverse are also needed.}

\noindent 子空间拓扑由 $Y\subseteq X$ 继承：$U\subseteq Y$ 开，当且仅当 $U=Y\cap V$，其中 $V$ 在 $X$ 中开。积拓扑由投影映射 $\pi_1:X\times Y\to X$、$\pi_2:X\times Y\to Y$ 的原像生成，基本开集是 $U\times V$。不交并拓扑把每个分支分别保留其拓扑。
\EN{A subspace inherits the topology $U=Y\cap V$ from the ambient space. The product topology is generated by the projections and has basis $U\times V$. A disjoint union keeps the topology of each component separately.}
\EX{$[0,1]$ 是 $\mathbb R$ 的子空间；集合 $[0,1)$ 在 $[0,1]$ 中开，因为它等于 $[0,1]\cap(-1,1)$，但在 $\mathbb R$ 中不是开集。}{The interval $[0,1]$ is a subspace of the real line. The set $[0,1)$ is open in $[0,1]$ because it is $[0,1]\cap(-1,1)$, although it is not open in $\mathbb R$.}

\section{同胚、紧性与连续像 | Homeomorphism and compactness}
\noindent 同胚是连续双射且逆映射连续的映射；同胚空间具有完全相同的拓扑性质。紧空间的开覆盖有有限子覆盖。对度量空间，紧性等价于序列紧性；在 $\mathbb R^n$ 中还等价于闭且有界。连续映射把紧集映为紧集，因此 Hausdorff 空间中的紧集闭。
\EN{A homeomorphism is a continuous bijection with a continuous inverse. Compactness means that every open cover has a finite subcover. In metric spaces it is equivalent to sequential compactness, and in $\mathbb R^n$ it is equivalent to being closed and bounded. Continuous images of compact sets are compact; compact subsets of Hausdorff spaces are closed.}
\EX{闭区间 $[0,1]$ 紧，而开区间 $(0,1)$ 不紧：开覆盖 $\{(1/n,1):n\ge2\}$ 覆盖 $(0,1)$，但没有有限子覆盖。连续函数 $f(x)=x$ 在 $[0,1]$ 上取得最大值和最小值。}{The closed interval $[0,1]$ is compact, while $(0,1)$ is not: the cover $\{(1/n,1):n\ge2\}$ has no finite subcover. The continuous function $f(x)=x$ attains its extrema on $[0,1]$.}

\noindent\textbf{Hausdorff 性。} 空间 Hausdorff，若任意两个不同点有不相交邻域。Hausdorff 空间中的序列极限唯一；例如通常的 $\mathbb R$ 是 Hausdorff，而平凡拓扑在含两个以上点时不是 Hausdorff。第一可数、第二可数和可分性是控制序列与基大小的条件；度量空间都是 Hausdorff 和第一可数。
\EN{A space is Hausdorff if distinct points have disjoint neighborhoods. Limits of sequences are unique in a Hausdorff space. The real line is Hausdorff, while a nontrivial trivial topology is not. First countability, second countability, and separability control the size of local or global bases; metric spaces are Hausdorff and first countable.}

\section{连通性与路径连通 | Connectedness and path connectedness}
\noindent 空间 $X$ 连通，若不能写成两个不交、非空、相对开集的并。等价地，$X$ 中既开又闭的集合只有 $\varnothing,X$。路径连通要求任意两点间存在连续路径 $\gamma:[0,1]\to X$；路径连通蕴含连通，但连通不一定路径连通。连续映射把连通集映为连通集。
\EN{A space is connected if it cannot be split into two disjoint nonempty open subsets. Equivalently, its only clopen sets are the empty set and the whole space. Path connectedness requires a continuous path between every two points. Path connectedness implies connectedness, but not conversely. Continuous images of connected spaces are connected.}
\EX{$[0,1]$ 是路径连通的；$(-1,0)\cup(0,1)$ 不连通，两个区间是分离。一个经典反例是拓扑学家的正弦曲线闭包，它连通但不是路径连通。}{The interval $[0,1]$ is path connected. The union $(-1,0)\cup(0,1)$ is disconnected. The closure of the topologist's sine curve is a standard connected but not path-connected example.}

\section{商拓扑与识别 | Quotient topology}
\noindent 给定满射 $q:X\to Y$，商拓扑定义为：$U\subseteq Y$ 开，当且仅当 $q^{-1}(U)$ 在 $X$ 中开。它把 $X$ 中被识别的点压成 $Y$ 中一个点。若 $f:X\to Z$ 连续且在每个纤维 $q^{-1}(y)$ 上取同一个值，则存在唯一映射 $\bar f:Y\to Z$ 使 $f=\bar f\circ q$；在商拓扑下 $\bar f$ 连续。
\EN{For a surjection $q:X\to Y$, the quotient topology declares $U$ open when $q^{-1}(U)$ is open in $X$. It identifies points in the same fiber. A continuous map constant on every fiber factors uniquely through the quotient, and the induced map is continuous.}
\EX{把区间 $[0,1]$ 的两个端点识别，得到一个圆周。把正方形的相对边按同向识别得到环面；把一对边反向识别则改变曲面的可定向性。}{Identifying the endpoints of $[0,1]$ produces a circle. Identifying opposite sides of a square in the same orientation produces a torus; reversing one identification changes orientability.}

\section{多边形粘合与曲面 | Polygon gluing and surfaces}
\noindent 多边形粘合时，先记录每条边的字母和方向，再判断顶点、边和面在商空间中的数量。一个紧致曲面的 Euler 特征为
\[
\chi=\#\{\text{vertices}\}-\#\{\text{edges}\}+\#\{\text{faces}\}.
\]
可定向闭曲面亏格 $g$ 的公式为 $\chi=2-2g$；非可定向闭曲面由 $k$ 个射影平面连通和得到，公式为 $\chi=2-k$。有 $n$ 个边界分支时分别减去 $n$。
\EN{For a polygon gluing, record edge labels and orientations, then count the resulting vertices, edges, and faces. The Euler characteristic is vertices minus edges plus faces. A closed orientable surface of genus $g$ has $\chi=2-2g$; a closed nonorientable surface with $k$ crosscaps has $\chi=2-k$. Each boundary component decreases the characteristic by one.}
\EX{正方形作为一个面，四个顶点全部识别为一个顶点，两对边识别为两条边，所以环面的 $\chi=1-2+1=0$。球面的四面体剖分有 $V=4,E=6,F=4$，所以 $\chi=2$。}{For the torus obtained from a square, all four vertices become one, the two edge pairs give two edges, and there is one face. Thus $\chi=1-2+1=0$. A tetrahedral decomposition of the sphere has $V=4,E=6,F=4$, giving $\chi=2$.}

\noindent 曲面可定向，直观上表示可以在整个曲面一致选择“正面”；在多边形图中，若边的识别全部保持方向，可得到可定向曲面，出现反向识别可能导致 Möbius 带、Klein 瓶或射影平面。分类定理说，紧致连通无边界曲面由球面、带 $g$ 个把手的可定向曲面、或带 $k$ 个交叉帽的非可定向曲面确定；实际分类还需检查可定向性。
\EN{A surface is orientable when a consistent choice of local orientation is possible. In a polygon scheme, orientation-reversing identifications can produce a Mobius band, a Klein bottle, or a projective plane. The classification theorem says that a compact connected boundaryless surface is a sphere, an orientable surface with handles, or a nonorientable surface with crosscaps; orientability must be checked separately.}

\section{同伦、变形收缩与 CW 复形 | Homotopy and CW complexes}
\noindent 映射 $f,g:X\to Y$ 同伦，若存在连续 $H:X\times[0,1]\to Y$，满足 $H(x,0)=f(x)$、$H(x,1)=g(x)$。若两个空间之间的映射复合后分别同伦于恒等映射，则它们同伦等价。同伦比同胚弱：同胚要求可逆的连续结构，同伦等价允许连续压缩。空间可缩，若它同伦等价于一点。
\EN{Maps are homotopic if a continuous map $H:X\times[0,1]\to Y$ deforms one into the other. Two spaces are homotopy equivalent when maps in both directions compose to maps homotopic to the identities. Homotopy equivalence is weaker than homeomorphism. A space is contractible if it is homotopy equivalent to a point.}
\EX{$\mathbb R^2$ 可通过 $H(x,t)=(1-t)x$ 收缩到原点，所以可缩；圆周不能这样收缩。圆柱和 Möbius 带都可变形收缩到圆周，因此同伦等价，但它们不一定同胚。}{The plane contracts to the origin by $H(x,t)=(1-t)x$, so it is contractible. A circle cannot be contracted in the same way. A cylinder and a Mobius band both deformation retract to a circle, so they are homotopy equivalent even though their global topology differs.}

\noindent CW 复形从离散的 $0$-胞腔开始，逐步粘合 $1$-胞腔、$2$-胞腔及更高维闭球。有限 CW 复形的 Euler 特征是胞腔数的交错和 $\chi=\sum_k(-1)^k c_k$，且对紧致 CW 复形是同伦不变量。有限连通图作为一维 CW 复形有 $\chi=V-E$；树可缩且 $\chi=1$，有环的图通常不具备可缩性。
\EN{A CW complex is built from discrete zero-cells by attaching one-cells, two-cells, and higher-dimensional disks. A finite CW complex has Euler characteristic $\chi=\sum_k(-1)^kc_k$, and for compact CW complexes it is a homotopy invariant. A finite connected graph has $\chi=V-E$; a tree is contractible and has characteristic one.}
\EX{圆周可看成一个 $0$-胞腔和一个把两端都粘到该点的 $1$-胞腔，因此 $\chi=1-1=0$；球面可看成一个 $0$-胞腔和一个 $2$-胞腔，故 $\chi=1+1=2$。}{A circle has one zero-cell and one one-cell, so $\chi=1-1=0$. A sphere can be built from one zero-cell and one two-cell, so $\chi=1+1=2$.}

\section{拓扑不变量与面试计算 | Invariants and oral calculations}
\noindent 拓扑不变量是同胚或同伦等价下不变的量。若两个空间的不变量不同，它们不可能同胚；若使用的是同伦不变量，不同值还能排除同伦等价。Euler 特征很容易计算，但不是完全不变量：不同空间可能有相同 Euler 特征，因此必要时需要基本群或同调群。
\EN{A topological invariant is unchanged under homeomorphism or, for a homotopy invariant, under homotopy equivalence. Different invariant values rule out the corresponding equivalence. Euler characteristic is easy to compute but is not a complete invariant; different spaces may have the same value, so fundamental groups or homology may be needed.}
\EX{球面和一个点的 Euler 特征分别为 $2$ 和 $1$，所以二者不同伦等价；圆环面和某些其他空间都可能有 $\chi=0$，仅凭 Euler 特征不能区分它们。}{The sphere and a point have Euler characteristics two and one, so they are not homotopy equivalent. A torus and other spaces may both have characteristic zero, so Euler characteristic alone may not distinguish them.}

\ASK{紧性为什么重要？}{Why is compactness important?}{它把无限开覆盖或无限序列转化为有限子覆盖或收敛子列，从而保证连续函数取得极值，并使很多存在性证明成立。}{Compactness turns an infinite cover into a finite subcover or gives a convergent subsequence. It guarantees extrema for continuous functions and supports many existence proofs.}
\ASK{同胚和同伦等价有什么区别？}{What is the difference between homeomorphism and homotopy equivalence?}{同胚要求连续双射且逆连续，保留完整拓扑结构；同伦等价只要求复合映射同伦于恒等映射，允许压缩可缩部分，因此更弱。}{A homeomorphism is a continuous bijection with continuous inverse and preserves the full topology. Homotopy equivalence only requires inverse maps up to homotopy and allows contractible parts to be compressed.}
\ASK{怎样由多边形粘合识别曲面？}{How do you identify a surface from a polygon gluing?}{先看边的方向判断可定向性，再数粘合后的顶点、边和面，计算 Euler 特征，最后用 $2-2g$ 或 $2-k$ 求亏格，并检查边界分支。}{Check orientations to determine orientability, count the resulting vertices, edges, and faces, compute Euler characteristic, and use $2-2g$ or $2-k$ to find the genus, including boundary components.}
\ASK{Euler 特征为什么不能单独分类所有空间？}{Why can Euler characteristic not classify all spaces?}{它只是一个整数，信息量有限；不同拓扑空间可能有相同 Euler 特征。它适合快速排除不可能的同伦类型，但需要基本群、同调等更强不变量完成区分。}{It is only one integer and loses information. Different spaces can have the same Euler characteristic. It is useful for ruling out possibilities, but stronger invariants such as fundamental groups or homology may be needed.}

\section{开集、闭集、内部与边界的计算 | Computing basic topological sets}
\noindent 对 $A\subseteq X$，内部 $A^\circ$ 是包含在 $A$ 中的最大开集；闭包 $\overline A$ 是包含 $A$ 的最小闭集；边界 $\partial A=\overline A\setminus A^\circ$。还可以用邻域描述：$x\in\overline A$ 当且仅当每个邻域都与 $A$ 相交；$x\in\partial A$ 当且仅当每个邻域都同时与 $A$ 和 $X\setminus A$ 相交。
\EN{For $A\subseteq X$, the interior is the largest open set contained in $A$, the closure is the smallest closed set containing $A$, and the boundary is closure minus interior. A point is in the closure when every neighborhood meets $A$; it is on the boundary when every neighborhood meets both $A$ and its complement.}
\EX{在 $\mathbb R$ 中，$A=(0,1]$ 的内部是 $(0,1)$，闭包是 $[0,1]$，边界是 $\{0,1\}$。有理数集 $\mathbb Q$ 的内部为空，闭包为 $\mathbb R$，所以边界也是 $\mathbb R$。}{In $\mathbb R$, the set $(0,1]$ has interior $(0,1)$, closure $[0,1]$, and boundary $\{0,1\}$. The rationals have empty interior and closure $\mathbb R$, so their boundary is all of $\mathbb R$.}

\noindent 常用恒等式是
\[
\overline{A\cup B}=\overline A\cup\overline B,
\qquad
(A\cup B)^c=A^c\cap B^c,
\qquad
\partial A=\overline A\cap\overline{X\setminus A}.
\]
闭包对有限并可交换，但对无限并只有 $\bigcup_i\overline{A_i}\subseteq\overline{\bigcup_iA_i}$，可能严格包含。这个区别是拓扑中有限性条件反复出现的原因。
\EN{Closure commutes with finite unions, but for an infinite union only $\bigcup_i\overline{A_i}\subseteq\overline{\bigcup_iA_i}$ is automatic. The inclusion may be strict, which explains many finite-versus-infinite distinctions in topology.}

\section{积拓扑与连续映射 | Product topology and maps}
\noindent 积空间 $X\times Y$ 的基本开集是 $U\times V$，其中 $U\subseteq X$、$V\subseteq Y$ 开。投影 $\pi_X(x,y)=x$、$\pi_Y(x,y)=y$ 连续。映射 $f:Z\to X\times Y$ 连续，当且仅当两个分量 $\pi_X\circ f$ 与 $\pi_Y\circ f$ 都连续；这称为积拓扑的泛性质。
\EN{The product topology has basis $U\times V$ with $U$ and $V$ open. The coordinate projections are continuous. A map into a product is continuous exactly when both coordinate maps are continuous.}
\EX{曲线 $\gamma(t)=(\cos t,\sin t)$ 从 $[0,2\pi]$ 到 $\mathbb R^2$ 连续，因为两个分量都是连续函数；其像是单位圆。}{The curve $\gamma(t)=(\cos t,\sin t)$ is continuous because both coordinate functions are continuous. Its image is the unit circle.}

\noindent 若 $X,Y$ 紧，则 $X\times Y$ 紧；更一般地，任意紧空间族的乘积紧，这是 Tychonoff 定理。课程中有限乘积可以用开覆盖和有限子覆盖证明。若 $X,Y$ Hausdorff，则 $X\times Y$ 也 Hausdorff，因为可以分别选不相交邻域并取它们的乘积。
\EN{A finite product of compact spaces is compact, and the general product theorem is Tychonoff's theorem. A finite product of Hausdorff spaces is Hausdorff because disjoint neighborhoods can be chosen in each coordinate.}

\section{连通性定理与证明思路 | Connectedness theorems}
\noindent 分离是 $X=U\cup V$，其中 $U,V$ 不交、非空且在 $X$ 中开。空间连通就是不存在分离。若 $A$ 连通且 $A\subseteq B\subseteq\overline A$，则 $B$ 连通；这说明连通集的闭包仍连通。连续像定理的证明使用反证：若 $f(A)$ 被分离，取原像会把 $A$ 分离。
\EN{A separation is a decomposition into two disjoint nonempty open sets. A space is connected when no such decomposition exists. A set between a connected set and its closure is connected, so closures of connected sets are connected. The continuous-image theorem follows by taking inverse images of a hypothetical separation.}
\EX{区间是连通的；如果区间被分成两个开闭非空部分，可以利用实数的上确界构造矛盾。连续函数的像仍是区间或单点，这也是介值定理的拓扑表达。}{Intervals are connected. A separation would contradict the least-upper-bound property of the real numbers. The continuous image of an interval is connected, hence an interval or a point; this is the topological form of the intermediate value theorem.}

\noindent 路径连通的定义是：任意 $x,y\in X$ 存在连续 $\gamma:[0,1]\to X$，满足 $\gamma(0)=x,\gamma(1)=y$。路径连通蕴含连通，因为如果路径跨越两个分离部分，$\gamma^{-1}(U)$ 和 $\gamma^{-1}(V)$ 会分离 $[0,1]$。反向不成立；拓扑学家的正弦曲线闭包是典型反例。
\EN{A path between any two points gives path connectedness. It implies connectedness because the interval $[0,1]$ is connected and a path cannot cross a separation. The converse fails; the closure of the topologist's sine curve is a standard connected but non-path-connected example.}

\section{紧性的几个等价观点 | Several views of compactness}
\noindent 开覆盖定义是：任意开覆盖 $\{U_i\}_{i\in I}$ 满足 $X=\bigcup_iU_i$，都存在有限指标 $i_1,\ldots,i_m$ 使 $X=U_{i_1}\cup\cdots\cup U_{i_m}$。在度量空间中，紧性等价于每个序列有收敛子列；在一般拓扑空间中，不能未经条件把这两个定义混用。紧集的闭子集紧，连续像紧。
\EN{Compactness means that every open cover has a finite subcover. In metric spaces it is equivalent to sequential compactness, but this equivalence may fail in general topological spaces. Closed subsets of compact spaces are compact, and continuous images of compact spaces are compact.}
\EX{用开覆盖证明 $[0,1]$ 紧的核心是取所有有限子覆盖能覆盖的右端点的上确界，再用开覆盖中的一个集合继续向右延伸，得到矛盾。这个证明使用了实数的完备性。}{The open-cover proof of compactness for $[0,1]$ takes the supremum of endpoints reachable by finite subcovers. A member of the cover extends past that supremum, giving a contradiction. The proof uses completeness of the real numbers.}

\noindent 紧性与连续性的结合给出三个常用结果：连续实值函数在紧集上有界并取得极值；连续函数从紧集到 Hausdorff 空间若为双射，则逆映射连续；紧集在 Hausdorff 空间中闭。第二个结果可以用于证明许多“连续双射其实是同胚”的结论。
\EN{Compactness combined with continuity gives boundedness and extrema for real-valued functions. A continuous bijection from a compact space to a Hausdorff space is a homeomorphism because it is a closed map. Compact subsets of Hausdorff spaces are closed.}

\section{分离公理与可数性 | Separation and countability}
\noindent $T_1$ 空间要求单点集闭；Hausdorff（$T_2$）空间要求任意两不同点有不相交邻域；正则空间可把点和不含它的闭集分离；正规空间可把两个不交闭集分离。通常度量空间满足 Hausdorff、正则和正规性质。一个容易混淆的点是：$T_1$ 不保证两个点有不相交邻域，Hausdorff 才保证。
\EN{$T_1$ means singleton sets are closed. Hausdorff means distinct points have disjoint neighborhoods. Regularity separates a point from a closed set, and normality separates two disjoint closed sets. Metric spaces satisfy these separation properties. $T_1$ alone does not imply Hausdorffness.}
\EX{有限补拓扑在无限集合上是 $T_1$，因为单点补集开；但不是 Hausdorff，因为两个非空开集总相交。它说明“点可以被闭集区分”弱于“点有不相交邻域”。}{The cofinite topology on an infinite set is $T_1$ because singleton complements are open, but it is not Hausdorff because any two nonempty open sets intersect.}

\noindent 第一可数表示每个点有可数邻域基；第二可数表示全空间有可数拓扑基；可分表示有可数稠密子集。第二可数空间可分；可度量空间若有可数基，常可用有理半径球描述。SEU 讲义强调这些条件，因为它们决定序列能否反映闭包、紧性和连续性。
\EN{First countability gives a countable local basis at each point. Second countability gives a countable basis for the whole topology. Separability means the existence of a countable dense subset. Second countability implies separability, and these assumptions make sequential arguments more reliable.}

\section{商拓扑的详细例子 | Detailed quotient examples}
\noindent 商拓扑的关键是“先取原像检查开性”。设 $q:[0,1]\to S^1$，$q(t)=e^{2\pi it}$，它把 $0$ 和 $1$ 识别。一个集合 $U\subseteq S^1$ 开，当且仅当 $q^{-1}(U)$ 在 $[0,1]$ 的子空间拓扑中开。端点的识别使圆周在局部仍像区间，但全局首尾相连。
\EN{For the quotient $q:[0,1]\to S^1$, $q(t)=e^{2\pi it}$, the endpoints are identified. A set in the circle is open exactly when its inverse image is open in the interval. The quotient is locally interval-like but globally has its ends joined.}

\noindent 商映射的泛性质可这样使用：若 $f:[0,1]\to Z$ 连续并满足 $f(0)=f(1)$，则它在被识别的点上取同值，所以存在唯一 $\bar f:S^1\to Z$ 使 $f=\bar f\circ q$。例如 $f(t)=(\cos2\pi t,\sin2\pi t)$ 下降为圆周到自身的恒等参数化。
\EN{If a continuous map on the interval satisfies $f(0)=f(1)$, it descends uniquely to a continuous map from the quotient circle. This universal property is often more useful than drawing the quotient.}

\section{曲面分类的计算流程 | A surface-classification workflow}
\noindent 对一个多边形粘合图，按以下顺序计算：
\begin{enumerate}
\item 记录每条边的配对和箭头方向；
\item 沿配对关系追踪顶点，确定商空间中顶点数 $V$；
\item 确定边数 $E$ 和面数 $F$；
\item 计算 $\chi=V-E+F$；
\item 由边方向判断可定向性，再使用闭曲面公式 $\chi=2-2g$ 或 $\chi=2-k$。
\end{enumerate}
\EN{For a polygon gluing, record paired edges and orientations, track vertices to find $V$, count edges $E$ and faces $F$, compute $\chi=V-E+F$, and then determine orientability before applying the genus formula.}

\EX{正方形相对边同向识别的环面：一个面 $F=1$，两对边给 $E=2$，四个顶点全部成为一个顶点 $V=1$，所以 $\chi=0$。它可定向，故 $0=2-2g$，得到 $g=1$。}{For the torus, one face remains, the two edge pairs give two edges, and all four vertices become one vertex. Thus $\chi=0$. The surface is orientable, so $0=2-2g$ gives genus one.}

\EX{射影平面可以由一个圆盘把边界上的对径点识别得到，也可以由正方形两对边反向粘合得到。它的胞腔分解有一个顶点、一条边和一个面，故 $\chi=1-1+1=1$；它不可定向，对应 $k=1$。}{The projective plane can be obtained from a disk by identifying opposite boundary points, or from a square with orientation-reversing edge identifications. Its cell structure has one vertex, one edge, and one face, so $\chi=1$. It is nonorientable and has one crosscap.}

\section{同伦与 Euler 特征的证明思路 | Homotopy and Euler characteristic}
\noindent 若 $X$ 变形收缩到 $A$，则包含映射 $i:A\hookrightarrow X$ 和收缩映射 $r:X\to A$ 满足 $r\circ i=\operatorname{id}_A$，且 $i\circ r\simeq\operatorname{id}_X$。因此 $X$ 与 $A$ 同伦等价。研究复杂空间时，可以先找一个更简单的变形收缩子空间。
\EN{If $X$ deformation retracts to $A$, the inclusion and retraction compose to the identity on $A$, while the other composition is homotopic to the identity on $X$. Hence $X$ and $A$ are homotopy equivalent. A complicated space can often be replaced by a simpler retract.}
\EX{圆柱 $S^1\times[0,1]$ 通过 $H((z,s),t)=(z,(1-t)s)$ 变形收缩到底圆周 $S^1\times\{0\}$。所以圆柱和圆周同伦等价，Euler 特征都为 $0$。}{The cylinder deformation retracts to its bottom circle by $H((z,s),t)=(z,(1-t)s)$. Thus the cylinder and circle have the same homotopy type and Euler characteristic zero.}

\noindent 有限 CW 复形的 Euler 特征 $\chi=\sum_k(-1)^kc_k$ 与具体剖分无关，而是同伦不变量。图的情形是 $\chi=V-E$；连通图若取一棵生成树，树有 $V-1$ 条边，每增加一条不在树中的边就增加一个独立环，因此独立环的数量为 $E-V+1=1-\chi$。
\EN{For a finite CW complex, the alternating cell count is independent of the chosen decomposition and is a homotopy invariant. For a connected graph, a spanning tree has $V-1$ edges; each remaining edge creates one independent cycle. Thus the number of independent cycles is $E-V+1=1-\chi$.}
\EX{一个有 $5$ 个顶点、$6$ 条边的连通图有 $6-5+1=2$ 个独立环，Euler 特征为 $-1$。它不能是树，也不能收缩到一点。}{A connected graph with five vertices and six edges has $6-5+1=2$ independent cycles and Euler characteristic $-1$. It is not a tree and cannot contract to a point.}

\section{全书面试复习的最短主线 | The shortest complete oral line}
\noindent 可以把整门拓扑学压缩成以下逻辑：拓扑空间由开集决定；开集定义连续性；连续映射保持连通性和紧性的相应结论；商拓扑把点识别起来并构造圆周、环面等空间；同伦允许压缩可缩部分；CW 复形提供计数方法；Euler 特征用顶点、边和面区分部分空间。每说一个概念，都准备一个具体空间和一个不满足条件的反例。
\EN{The complete course can be summarized as follows: open sets define topology and continuity; continuous maps preserve connectedness and compactness in the appropriate direction; quotient topology constructs circles and surfaces by identifying points; homotopy permits contraction; CW complexes provide a counting model; Euler characteristic distinguishes some spaces. For every concept, prepare one example and one counterexample.}

\ASK{连续双射为什么不一定是同胚？}{Why is a continuous bijection not always a homeomorphism?}{还需要逆映射连续。若定义域紧、陪域 Hausdorff，则连续双射自动是同胚；没有这些条件就不能直接推出。}{The inverse must also be continuous. A continuous bijection from a compact space to a Hausdorff space is a homeomorphism, but without these hypotheses the conclusion may fail.}
\ASK{为什么 $\mathbb R$ 连通但 $\mathbb Q$ 也看起来没有断开？}{Why is the real line connected but the rationals are not?}{$\mathbb R$ 的完备性保证任何分离都会违反上确界性质；$\mathbb Q$ 缺少无理数极限，例如可以用小于 $\sqrt2$ 和大于 $\sqrt2$ 的有理数把它分成两部分。}{Completeness of $\mathbb R$ rules out a separation using the least-upper-bound property. The rationals can be separated by rational numbers below and above $\sqrt2$, because the missing irrational point creates a gap.}
\ASK{为什么 Euler 特征是拓扑不变量？}{Why is Euler characteristic a topological invariant?}{对有限 CW 复形，细分会同时增加不同维度的胞腔，新增数量的交错和相互抵消；更一般地它在同伦等价下保持不变。课程中可先使用定理，再理解图和曲面的计算。}{For finite CW complexes, subdivision adds cells in different dimensions whose alternating contributions cancel. More generally Euler characteristic is homotopy invariant. In an oral exam, state the theorem and then demonstrate it on graphs and surfaces.}

\section{从零理解“拓扑”到底保留什么 | What topology keeps}
\noindent 拓扑学先忘掉长度和角度，只保留哪些点彼此靠近、哪些集合可以视为一个开区域。一个拓扑空间是集合 $X$ 加上一族开集 $\mathcal T$。因此“点 $x$ 附近发生什么”要用包含 $x$ 的开集描述，而不是必须使用半径。邻域可以很大；只要它包含某个开集 $U$ 且 $x\in U$，就可以称为 $x$ 的邻域。
\EN{Topology forgets lengths and angles and keeps the structure of neighborhoods. A topological space consists of a set and a family of open sets. Local statements are described by open neighborhoods rather than by a particular radius or metric.}

\EX{在通常的实数拓扑中，点 $1$ 的一个邻域可以是 $(0,2)$，也可以是 $(-100,100)$；它们大小不同，但都包含一个围绕 $1$ 的开区间。在离散拓扑中，单点集 $\{1\}$ 本身就是开集，所以每个点都可以被完全分离。}{In the usual topology on the real line, both $(0,2)$ and $(-100,100)$ are neighborhoods of one. In the discrete topology, the singleton containing one is open, so every point can be isolated.}

\noindent 连续映射的核心不是“距离变化很小”，而是\textbf{开集的原像仍然开}：若 $f:X\to Y$ 连续，$V\subseteq Y$ 开，则 $f^{-1}(V)$ 在 $X$ 中开。这个定义可以同时处理实函数、曲线、投影和商映射。若 $f$ 是双射且 $f^{-1}$ 也连续，才称为同胚；同胚保留开闭、紧性、连通性等拓扑性质。
\EN{Continuity means that inverse images of open sets are open. This definition handles real functions, curves, projections, and quotient maps in the same way. A bijective continuous map is a homeomorphism only when its inverse is also continuous.}

\EX{投影 $\pi_1:X\times Y\to X$ 满足 $\pi_1^{-1}(U)=U\times Y$。积拓扑正是让这些投影连续的最弱拓扑。对于 $f:Z\to X\times Y$，只要两个分量 $\pi_1\circ f$ 和 $\pi_2\circ f$ 连续，$f$ 就连续。}{For the projection $\pi_1:X\times Y\to X$, the inverse image of $U$ is $U\times Y$. The product topology is the weakest topology making the projections continuous. A map into a product is continuous exactly when both coordinate maps are continuous.}

\section{条件不能省：四个核心定理的使用范围 | The hypotheses matter}
\noindent 拓扑面试中最容易出错的是把定理的结论说得过强。应当同时记住以下条件：
\begin{enumerate}
\item 在任意拓扑空间中，连续像保持紧性和连通性；
\item 在 Hausdorff 空间中，紧子集是闭的；
\item 在紧空间中，连续实值函数取得最大值和最小值；
\item “闭且有界等价于紧”只在 $\mathbb R^n$ 等特定空间成立，不是任意拓扑空间的定义。
\end{enumerate}
\EN{The most common mistake is to omit hypotheses. Continuous images preserve compactness and connectedness. Compact subsets of Hausdorff spaces are closed. Continuous real-valued functions attain extrema on compact spaces. The statement closed and bounded if and only if compact is special to Euclidean spaces, not a general definition of compactness.}

\EX{在 $\mathbb R$ 中，$[0,1]$ 闭且有界，所以紧；在无限维函数空间中，闭单位球通常不紧。另一方面，$(0,1)$ 有界但不闭且不紧。这个比较说明“有界”只依赖距离，而紧性是更强的开覆盖性质。}{In the real line, $[0,1]$ is closed and bounded and hence compact. In many infinite-dimensional function spaces, the closed unit ball is not compact. The interval $(0,1)$ is bounded but not compact. Boundedness alone is therefore much weaker than compactness.}

\noindent \textbf{证明时的三条常用路线。} 证明连续像连通，假设像有分离，再取原像；证明连续像紧，把目标空间的开覆盖拉回定义域并取有限子覆盖；证明紧集上的连续函数取极值，先把函数像看成紧的，再使用实数完备性。回答时先说路线，再写关键集合关系，能避免只背结论。
\EN{Three standard proof patterns are useful. For a connected image, pull back a hypothetical separation. For a compact image, pull back an open cover and take a finite subcover. For extrema, show that the image is compact and use completeness of the real line.}

\section{商空间的逐步构造 | Constructing quotient spaces step by step}
\noindent 商空间的实际操作可以按四步进行：
\begin{enumerate}
\item 写出原空间 $X$ 和等价关系，说明哪些点要识别；
\item 写出商映射 $q:X\to X/{\sim}$；
\item 用 $U$ 在商空间开 $\Longleftrightarrow q^{-1}(U)$ 在 $X$ 中开来检查开集；
\item 若要证明映射从商空间下降，验证它在每个等价类上取同一个值。
\end{enumerate}
\EN{To construct a quotient, specify the original space and the equivalence relation, define the quotient map, test openness by inverse images, and check that a map is constant on equivalence classes before descending it to the quotient.}

\EX{令 $q:[0,1]\to S^1$，$q(t)=e^{2\pi it}$。只有 $0$ 和 $1$ 被识别。若 $f:[0,1]\to Z$ 连续且 $f(0)=f(1)$，定义 $\bar f(q(t))=f(t)$。等价关系保证这个定义不依赖代表元；商拓扑的泛性质保证 $\bar f$ 连续。}{Let $q:[0,1]\to S^1$ be $q(t)=e^{2\pi it}$, identifying only zero and one. If $f:[0,1]\to Z$ is continuous and $f(0)=f(1)$, define $\bar f(q(t))=f(t)$. The endpoint condition makes this well-defined, and the quotient property makes the induced map continuous.}

\noindent 对正方形还要记录边的方向。相对边按 $a,b,a^{-1},b^{-1}$ 粘合得到环面；若出现同一字母的边按相同方向粘合，通常会产生不可定向曲面。数顶点时不能只数原多边形的角，而要按照边的识别关系追踪它们在商空间中的等价类。
\EN{For a square, record both labels and orientations. The word $aba^{-1}b^{-1}$ gives the torus. Same-direction identifications can produce nonorientable surfaces. When counting vertices, count equivalence classes after gluing, rather than the original corners.}

\EX{Möbius 带可由矩形的一对相对边反向识别得到。它有一个边界分支，不能定向；它可以变形收缩到中心圆周，因此其 Euler 特征为 $0$。这与环面同样的 Euler 特征说明 Euler 特征不能单独判断可定向性。}{A Mobius band is obtained by reversing one pair of opposite sides of a rectangle. It has one boundary component and is nonorientable. It deformation retracts to its central circle, so its Euler characteristic is zero. The same characteristic as the torus shows that Euler characteristic alone does not determine orientability.}

\section{同伦计算：如何真正写出一个变形 | How to write a homotopy}
\noindent 同伦 $H:X\times[0,1]\to Y$ 必须同时满足三件事：对每个 $t$ 是一个映射，对 $(x,t)$ 联合连续，并且 $H(x,0)=f(x)$、$H(x,1)=g(x)$。变形收缩还要求子空间 $A$ 中的点在整个过程中保持不动，即 $H(a,t)=a$。最后一条经常被遗漏。
\EN{A homotopy $H:X\times[0,1]\to Y$ must be jointly continuous and satisfy the two endpoint conditions. A deformation retraction additionally fixes every point of the target subspace throughout the deformation.}

\EX{圆柱 $X=S^1\times[0,1]$ 到底圆周 $A=S^1\times\{0\}$ 的变形收缩为
\[
H((z,s),t)=(z,(1-t)s).
\]
当 $t=0$ 时是恒等映射，当 $t=1$ 时落在 $A$ 中；若 $s=0$，始终有 $H((z,0),t)=(z,0)$，所以确实固定 $A$。}{The cylinder deformation retracts to its bottom circle by $H((z,s),t)=(z,(1-t)s)$. At time zero it is the identity, at time one it lies in the bottom circle, and points already in the bottom circle stay fixed.}

\noindent 变形收缩常被用来简化计算：树收缩到一点，圆柱收缩到圆周，带孔平面不能直接收缩到一点，因为孔造成了不能消去的环路。课程中若尚未定义基本群，可以先把这句话理解为“存在闭曲线障碍”；若继续学习代数拓扑，再用基本群把障碍写成代数对象。
\EN{Deformation retractions simplify spaces: a tree contracts to a point and a cylinder contracts to a circle. A punctured plane cannot contract to a point because a loop around the hole cannot be removed. If fundamental groups have not yet been introduced, this can be understood as a persistent loop obstruction.}

\section{曲面和图的计算练习 | Worked calculations}
\EX{一个连通图有 $V=6$ 个顶点、$E=7$ 条边。它作为一维 CW 复形的 Euler 特征是 $\chi=6-7=-1$，独立环的数量是 $E-V+1=2$。若从一棵生成树开始，树有 $5$ 条边，剩下的两条边各增加一个环。}{A connected graph with six vertices and seven edges has Euler characteristic $6-7=-1$ and two independent cycles. A spanning tree has five edges, and the two remaining edges create the two cycles.}

\EX{若一个闭可定向曲面的胞腔计数给出 $V=1,E=6,F=2$，则 $\chi=1-6+2=-3$。由 $2-2g=-3$ 得 $g=5/2$，不是整数，因此这些计数或可定向闭曲面的假设有错误。这个检查可以迅速发现粘合计算中的漏数。}{If a proposed closed orientable surface has $V=1$, $E=6$, and $F=2$, then $\chi=-3$. The equation $2-2g=-3$ gives a noninteger genus, so the cell count or the orientability and closedness assumptions must be wrong. This is a quick consistency check.}

\ASK{为什么商拓扑要用原像定义开集？}{Why is the quotient topology defined using inverse images?}{因为原像运算自动保持并、交和补的关系，能够使商映射连续；同时它是使 $q$ 连续的最细拓扑。若直接凭图形判断开集，容易忽略端点或边界被识别后的邻域。}{Inverse images preserve unions, finite intersections, and complements, and this makes the quotient map continuous. The quotient topology is the finest topology with this property. Visual intuition alone can miss neighborhoods created by identifying boundary points.}
\ASK{怎样判断一个 Euler 特征计算是否可信？}{How can you check an Euler characteristic calculation?}{先确认空间是否是有限 CW 复形或有限剖分，再按识别关系数顶点和边，最后检查公式给出的亏格是否为合法整数，并核对是否有边界和是否可定向。}{First verify that a finite cell decomposition is being used. Count equivalence classes of vertices and edges, check whether the resulting genus is an allowed integer, and account for boundary components and orientability.}

\section*{来源与范围 | Source and scope}
本稿覆盖“拓扑学”文件夹中的 SEU《Topology》讲义全书：第 1 章一般拓扑学（度量空间、有限维空间、拓扑与基、子空间、积空间、连通性、紧性、分离性与可数性）以及第 2 章几何与代数拓扑基础（商拓扑、多边形粘合、曲面分类、同伦等价、CW 复形和 Euler 特征）。代数拓扑另有独立的第 1--5 章稿件；本稿为点集拓扑和曲面基础的课程复习，不把两份内容混作同一门课。
\EN{This review covers the complete SEU Topology notes: Chapter 1 on general topology and Chapter 2 on rudiments of geometric and algebraic topology. The separate algebraic-topology review covers simplicial homology, fundamental groups, and coverings. This document focuses on point-set topology and the foundations of surfaces, without treating the two courses as identical.}
\end{document}
